\documentclass[11pt,a4paper]{article}

\usepackage[utf8]{inputenc}
\usepackage[T1]{fontenc}
\usepackage[french]{babel}
\usepackage[margin=2cm]{geometry}
\usepackage{parskip}
\usepackage{amsmath,amssymb,amsthm}
\usepackage{mathtools}
\usepackage{booktabs}
\usepackage{enumitem}
\usepackage{listings}
\usepackage{xcolor}
\usepackage{hyperref}
\hypersetup{colorlinks=true,linkcolor=blue!70!black}

\theoremstyle{plain}
\newtheorem{theorem}{Théorème}[section]
\newtheorem{lemma}[theorem]{Lemme}
\newtheorem{proposition}[theorem]{Proposition}
\newtheorem{corollary}[theorem]{Corollaire}
\theoremstyle{definition}
\newtheorem{definition}[theorem]{Définition}
\newtheorem{hypothesis}[theorem]{Hypothèse}
\theoremstyle{remark}
\newtheorem*{remark}{Remarque}
\newtheorem{observation}[theorem]{Observation}

\newcommand{\R}{\mathbb{R}}
\newcommand{\N}{\mathbb{N}}
\newcommand{\E}{\mathbb{E}}
\newcommand{\Prob}{\mathbb{P}}
\newcommand{\Var}{\mathrm{Var}}
\newcommand{\Cov}{\mathrm{Cov}}
\newcommand{\ind}{\mathbf{1}}
\newcommand{\calL}{\mathcal{L}}
\newcommand{\calF}{\mathcal{F}}
\newcommand{\calS}{\mathcal{S}}
\newcommand{\bigO}{\mathcal{O}}

\lstset{basicstyle=\ttfamily\small,language=C,frame=single}

\title{\textbf{Preuve de Convergence pour le Modèle V2} \\[0.5em]
\Large Système de Buffers Indépendants (No Producer Contention) \\
MPMC Ring Buffer V2}
\author{Document Technique}
\date{Janvier 2025}

\begin{document}
\maketitle

\begin{abstract}
Ce document présente la preuve de convergence pour le modèle V2 du MPMC ring buffer, basé sur des \textbf{buffers per-CPU indépendants} sans contention producteur. Contrairement au modèle V1 (mean-field), le modèle V2 repose sur la \textbf{Loi des Grands Nombres} pour N processus i.i.d., ce qui simplifie la preuve et améliore la vitesse de convergence (exponentielle au lieu de polynomiale).
\end{abstract}

\tableofcontents
\newpage

%%%%%%%%%%%%%%%%%%%%%%%%%%%%%%%%%%%%%%%%%%%%%%%%%%%%%%%%%%%%%%%%%%%%%%%%%%%%%%%
\part{Le Modèle V2 : No Producer Contention}
%%%%%%%%%%%%%%%%%%%%%%%%%%%%%%%%%%%%%%%%%%%%%%%%%%%%%%%%%%%%%%%%%%%%%%%%%%%%%%%

\section{Motivation : Élimination de l'État Global}

\subsection{Problème du modèle V1 (Mean-Field)}

Dans le modèle V1, tous les producteurs accèdent à un \textbf{état global partagé} :

\begin{lstlisting}
// V1 : Etat global partage
atomic_t prod_head;      // Contention!
atomic_t ready_count;    // Contention!

// Tous les producteurs font:
start = atomic_fetch_add(&prod_head, 1);  // CAS global
\end{lstlisting}

Cela crée :
\begin{itemize}
    \item De la \textbf{contention} sur les opérations atomiques
    \item Des \textbf{boucles de retry} (non wait-free)
    \item Une \textbf{dégradation} avec le nombre de threads
\end{itemize}

\subsection{Solution V2 : Buffers Per-CPU Indépendants}

Le modèle V2 élimine la \textbf{contention entre producteurs} :

\begin{lstlisting}
// V2 : No producer contention (but minimal consumer shared state)
struct per_cpu_buffer buffers[NUM_CPUS];

// Chaque producteur ecrit UNIQUEMENT dans son buffer local:
cpu_id = raw_smp_processor_id();
h = &buffers[cpu_id];  // Notre buffer exclusif
cmpxchg(&h->state_gen, cur, next);  // CAS local seulement
\end{lstlisting}

\textbf{Conséquence fondamentale} : Les N buffers évoluent de manière \textbf{indépendante}.

\section{Description Formelle du Système V2}

\begin{definition}[Système V2 à N buffers indépendants]
On considère $N$ buffers, un par CPU. Chaque buffer $i \in \{1, \ldots, N\}$ évolue selon une chaîne de Markov \textbf{indépendante} avec états $\calS = \{E, F, R, D\}$ :
\begin{itemize}
    \item $E$ (Empty) : Buffer disponible pour écriture
    \item $F$ (Filling) : Buffer en cours d'écriture par le producteur local
    \item $R$ (Ready) : Données prêtes, en attente de consommation
    \item $D$ (Reading) : Buffer en cours de lecture par un consommateur
\end{itemize}
\end{definition}

\begin{definition}[Transitions du buffer $i$]
Le buffer $i$ effectue les transitions suivantes avec des taux \textbf{indépendants des autres buffers} :
\begin{align}
E \xrightarrow{\lambda_P} F & \quad \text{(réservation par producteur local)} \\
F \xrightarrow{\lambda_S} R & \quad \text{(submit)} \\
F \xrightarrow{\lambda_d} E & \quad \text{(discard)} \\
R \xrightarrow{\lambda_C} D & \quad \text{(claim par consommateur)} \\
D \xrightarrow{\mu} E & \quad \text{(release)}
\end{align}
\end{definition}

\begin{observation}[Indépendance --- Validité]
\textbf{Cruciale pour les producteurs} : Le taux de transition du buffer $i$ ne dépend \textbf{que de l'état de ce buffer} pour les opérations producteur.

\textbf{Pour les consommateurs} : L'implémentation actuelle utilise \texttt{\_Thread\_local} pour la rotation des partitions. Il n'y a \textbf{pas d'état global partagé} entre consommateurs au niveau de la queue.

Le modèle i.i.d. est donc une \textbf{bonne approximation} car :
\begin{itemize}
    \item Pas de contention sur des atomiques globaux
    \item Chaque consumer a sa propre rotation locale
    \item La seule interaction est via les CAS sur les slots individuels
\end{itemize}
\end{observation}

\section{Comparaison V1 vs V2}

\begin{table}[htbp]
\centering
\caption{Différences fondamentales entre les modèles V1 et V2.}
\label{tab:v1_v2_comparison}
\small
\begin{tabular}{lcc}
\toprule
\textbf{Aspect} & \textbf{V1 (Mean-Field)} & \textbf{V2 (Indépendant)} \\
\midrule
État global & $\text{prod\_head}, \text{ready\_count}$ & \textbf{Aucun} \\
Taux de transition & $r_j(\sigma) = N \cdot \tilde{r}_j(\rho^N)$ & $r_j^{(i)} = r_j(\sigma_i)$ \\
Dépendance & Tous les buffers interagissent & Buffers indépendants \\
Échangeabilité (H1) & Requise & \textbf{Non applicable} \\
Density-dependent (H3) & Requise & \textbf{Non applicable} \\
Mécanisme de preuve & Kurtz (mean-field limit) & LGN (i.i.d.) \\
Vitesse de convergence & $O(1/\sqrt{N})$ & $\mathbf{O(e^{-cN})}$ \\
Garantie lock-free & Non & \textbf{Oui (borné)} \\
\bottomrule
\end{tabular}
\end{table}

%%%%%%%%%%%%%%%%%%%%%%%%%%%%%%%%%%%%%%%%%%%%%%%%%%%%%%%%%%%%%%%%%%%%%%%%%%%%%%%
\part{La Preuve pour le Modèle V2}
%%%%%%%%%%%%%%%%%%%%%%%%%%%%%%%%%%%%%%%%%%%%%%%%%%%%%%%%%%%%%%%%%%%%%%%%%%%%%%%

\section{Structure de la Preuve}

La preuve V2 est \textbf{plus simple} que la preuve V1 car elle repose sur :
\begin{enumerate}
    \item La \textbf{Loi des Grands Nombres} pour des variables i.i.d.
    \item Les \textbf{équations de Kolmogorov} pour une chaîne de Markov unique
    \item L'inégalité de \textbf{Hoeffding} pour la concentration
\end{enumerate}

Aucune hypothèse d'échangeabilité (H1), de density-dependence (H3), ou de mélange rapide (H4) n'est nécessaire.

\section{Étape 1 : Dynamique d'un Buffer Individuel}

\begin{definition}[Chaîne de Markov d'un buffer]
Chaque buffer $i$ est une chaîne de Markov à temps continu sur $\calS = \{E, F, R, D\}$ avec matrice génératrice :
\[
Q = \begin{pmatrix}
-\lambda_P & \lambda_P & 0 & 0 \\
\lambda_d & -(\lambda_S + \lambda_d) & \lambda_S & 0 \\
0 & 0 & -\lambda_C & \lambda_C \\
\mu & 0 & 0 & -\mu
\end{pmatrix}
\]
\end{definition}

\begin{lemma}[Distribution stationnaire d'un buffer]
La distribution stationnaire $\pi = (\pi_E, \pi_F, \pi_R, \pi_D)$ est l'unique solution de $\pi Q = 0$ avec $\sum_s \pi_s = 1$.
\end{lemma}

\begin{proof}
Le système $\pi Q = 0$ donne :
\begin{align}
0 &= -\lambda_P \pi_E + \lambda_d \pi_F + \mu \pi_D \\
0 &= \lambda_P \pi_E - (\lambda_S + \lambda_d) \pi_F \\
0 &= \lambda_S \pi_F - \lambda_C \pi_R \\
0 &= \lambda_C \pi_R - \mu \pi_D
\end{align}

De (2) : $\pi_F = \frac{\lambda_P}{\lambda_S + \lambda_d} \pi_E$.

De (3) : $\pi_R = \frac{\lambda_S}{\lambda_C} \pi_F = \frac{\lambda_S \lambda_P}{\lambda_C(\lambda_S + \lambda_d)} \pi_E$.

De (4) : $\pi_D = \frac{\lambda_C}{\mu} \pi_R = \frac{\lambda_S \lambda_P}{\mu(\lambda_S + \lambda_d)} \pi_E$.

La normalisation $\pi_E + \pi_F + \pi_R + \pi_D = 1$ détermine $\pi_E$.
\end{proof}

\begin{definition}[Probabilités transitoires]
Pour un buffer dans l'état $s$ à $t=0$, la probabilité d'être dans l'état $s'$ au temps $t$ est :
\[
P_{ss'}(t) = (e^{Qt})_{ss'}
\]

La probabilité marginale à temps $t$ est :
\[
\pi_s(t) = \sum_{s'} \pi_{s'}(0) \cdot P_{s's}(t)
\]
\end{definition}

\begin{theorem}[Convergence vers l'équilibre]\label{thm:single_buffer_convergence}
Pour toute condition initiale, $\pi(t) \to \pi^*$ exponentiellement vite :
\[
\|\pi(t) - \pi^*\|_{TV} \leq C \cdot e^{-\gamma t}
\]
où $\gamma = \min\{|\text{Re}(\lambda)| : \lambda \text{ valeur propre non nulle de } Q\} > 0$.
\end{theorem}

\begin{proof}
$Q$ est une matrice génératrice irréductible sur un espace fini. Par le théorème de Perron-Frobenius, 0 est valeur propre simple et toutes les autres valeurs propres ont une partie réelle strictement négative.

La solution de $\frac{d\pi}{dt} = \pi Q$ est $\pi(t) = \pi(0) e^{Qt}$, qui converge vers $\pi^*$ à vitesse $e^{-\gamma t}$.
\end{proof}

\section{Étape 2 : Densités Empiriques comme Moyenne i.i.d.}

\begin{definition}[Densité empirique]
La densité empirique de l'état $s$ au temps $t$ est :
\[
\rho^N_s(t) = \frac{1}{N} \sum_{i=1}^{N} \ind_{\sigma_i(t) = s}
\]
où $\sigma_i(t) \in \calS$ est l'état du buffer $i$ au temps $t$.
\end{definition}

\begin{observation}[Structure i.i.d.]
Puisque les buffers sont \textbf{indépendants} et ont la \textbf{même loi} (même chaîne de Markov), les variables $\ind_{\sigma_i(t) = s}$ sont \textbf{i.i.d.} avec :
\[
\E[\ind_{\sigma_i(t) = s}] = \pi_s(t)
\]
\end{observation}

\begin{theorem}[Loi des Grands Nombres]\label{thm:lgn}
Pour tout $t \geq 0$ et tout $s \in \calS$ :
\[
\rho^N_s(t) \xrightarrow[N \to \infty]{p.s.} \pi_s(t)
\]
\end{theorem}

\begin{proof}
C'est la Loi des Grands Nombres classique pour des variables i.i.d. bornées (Bernoulli).

Soit $X_i = \ind_{\sigma_i(t) = s}$. Les $X_i$ sont i.i.d. avec $\E[X_i] = \pi_s(t)$.

Par la LGN forte :
\[
\frac{1}{N} \sum_{i=1}^N X_i \xrightarrow{p.s.} \E[X_1] = \pi_s(t)
\]
\end{proof}

\section{Étape 3 : Convergence Exponentielle (Hoeffding)}

\begin{theorem}[Concentration exponentielle]\label{thm:hoeffding}
Pour tout $\epsilon > 0$, $t \geq 0$, et $s \in \calS$ :
\[
\boxed{\Prob\left( |\rho^N_s(t) - \pi_s(t)| > \epsilon \right) \leq 2 \cdot e^{-2N\epsilon^2}}
\]
\end{theorem}

\begin{proof}
Les variables $X_i = \ind_{\sigma_i(t) = s}$ sont i.i.d. avec $X_i \in [0, 1]$.

Par l'inégalité de Hoeffding pour des variables bornées :
\[
\Prob\left( \left| \frac{1}{N}\sum_{i=1}^N X_i - \E[X_1] \right| > \epsilon \right) \leq 2 \exp\left( -\frac{2N\epsilon^2}{\sum_i (b_i - a_i)^2 / N^2} \right)
\]

Avec $a_i = 0$, $b_i = 1$ :
\[
\Prob\left( |\rho^N_s(t) - \pi_s(t)| > \epsilon \right) \leq 2 e^{-2N\epsilon^2}
\]
\end{proof}

\begin{corollary}[Comparaison avec V1]
\begin{itemize}
    \item \textbf{V1 (mean-field)} : $\Prob(|\rho^N - \rho| > \epsilon) \leq \frac{C}{N\epsilon^2}$ (polynomiale)
    \item \textbf{V2 (indépendant)} : $\Prob(|\rho^N - \pi| > \epsilon) \leq 2e^{-2N\epsilon^2}$ (exponentielle)
\end{itemize}

La convergence V2 est \textbf{exponentiellement plus rapide} !
\end{corollary}

\section{Étape 4 : L'ODE Limite}

\begin{theorem}[Équivalence avec l'ODE mean-field]\label{thm:ode_equivalence}
Les probabilités transitoires $\pi(t)$ satisfont la même ODE que la limite mean-field :
\[
\frac{d\pi}{dt} = F(\pi)
\]
où $F$ est le champ de vecteurs mean-field défini par :
\[
F(\pi) = \begin{pmatrix}
\mu \pi_D + \lambda_d \pi_F - \lambda_P \pi_E \\
\lambda_P \pi_E - (\lambda_S + \lambda_d) \pi_F \\
\lambda_S \pi_F - \lambda_C \pi_R \\
\lambda_C \pi_R - \mu \pi_D
\end{pmatrix}
\]
\end{theorem}

\begin{proof}
L'équation de Kolmogorov forward (master equation) pour une chaîne de Markov est :
\[
\frac{d\pi}{dt} = \pi Q
\]

En développant $\pi Q$ composante par composante, on obtient exactement $F(\pi)$.
\end{proof}

\begin{remark}[Même limite, mécanisme différent]
Les modèles V1 et V2 ont la \textbf{même ODE limite} $\dot{\rho} = F(\rho)$.

Mais le \textbf{mécanisme de convergence} est différent :
\begin{itemize}
    \item V1 : L'interaction mean-field concentre les densités autour de la solution de l'ODE
    \item V2 : L'indépendance + LGN concentre les densités autour des probabilités marginales, qui satisfont l'ODE
\end{itemize}
\end{remark}

\section{Étape 5 : Convergence Uniforme en Temps}

\begin{theorem}[Convergence uniforme]\label{thm:uniform_convergence}
Pour tout $T > 0$ et $\epsilon > 0$ :
\[
\boxed{\Prob\left( \sup_{0 \leq t \leq T} \|\rho^N(t) - \pi(t)\| > \epsilon \right) \leq C_T \cdot e^{-cN\epsilon^2}}
\]
où $C_T$ dépend de $T$ mais $c > 0$ est une constante universelle.
\end{theorem}

\begin{proof}
\textbf{Étape 1 : Discrétisation temporelle.}

On découpe $[0, T]$ en $M = \lceil T/\delta \rceil$ intervalles de longueur $\delta$.

Soit $t_k = k\delta$ pour $k = 0, \ldots, M$.

\textbf{Étape 2 : Union bound sur les points de grille.}

Par Hoeffding (Théorème~\ref{thm:hoeffding}) et union bound :
\[
\Prob\left( \exists k : \|\rho^N(t_k) - \pi(t_k)\| > \epsilon/2 \right) \leq M \cdot 8 e^{-N\epsilon^2/2}
\]

(le facteur 8 vient de $2 \times 4$ états)

\textbf{Étape 3 : Contrôle entre les points de grille.}

Entre $t_k$ et $t_{k+1}$, la variation de $\rho^N$ est bornée par le nombre de transitions.

Par couplage, si $\|\rho^N(t_k) - \pi(t_k)\| \leq \epsilon/2$, alors pour $\delta$ assez petit :
\[
\|\rho^N(t) - \pi(t)\| \leq \epsilon/2 + O(\delta) \leq \epsilon
\]

\textbf{Étape 4 : Conclusion.}

En choisissant $\delta = O(\epsilon)$ et $M = O(T/\epsilon)$, on obtient :
\[
\Prob\left( \sup_{t \leq T} \|\rho^N - \pi\| > \epsilon \right) \leq \frac{CT}{\epsilon} \cdot e^{-cN\epsilon^2}
\]
\end{proof}

%%%%%%%%%%%%%%%%%%%%%%%%%%%%%%%%%%%%%%%%%%%%%%%%%%%%%%%%%%%%%%%%%%%%%%%%%%%%%%%
\part{Théorème Central Limite}
%%%%%%%%%%%%%%%%%%%%%%%%%%%%%%%%%%%%%%%%%%%%%%%%%%%%%%%%%%%%%%%%%%%%%%%%%%%%%%%

\section{Fluctuations Gaussiennes}

\begin{definition}[Processus de fluctuations]
\[
\xi^N(t) = \sqrt{N} \left( \rho^N(t) - \pi(t) \right)
\]
\end{definition}

\begin{theorem}[TCL pour le modèle V2]\label{thm:clt_v2}
Le processus $\xi^N(t)$ converge en loi vers un processus gaussien centré avec covariance :
\[
\Gamma_t(s, s') = \pi_s(t) \ind_{s=s'} - \pi_s(t) \pi_{s'}(t)
\]
\end{theorem}

\begin{proof}
C'est le TCL classique pour des sommes i.i.d.

Pour chaque $t$ fixé, les $X_i = \ind_{\sigma_i(t) = s}$ sont i.i.d. Bernoulli($\pi_s(t)$).

Par le TCL :
\[
\sqrt{N}(\rho^N_s(t) - \pi_s(t)) = \frac{1}{\sqrt{N}} \sum_i (X_i - \pi_s(t)) \xrightarrow{\calL} \mathcal{N}(0, \pi_s(t)(1 - \pi_s(t)))
\]

Le vecteur joint $\xi^N(t)$ converge vers un vecteur gaussien avec covariance multinomiale.
\end{proof}

\begin{remark}[Différence avec V1]
Dans V1, les fluctuations suivent un processus d'Ornstein-Uhlenbeck (corrélées dans le temps via l'interaction mean-field).

Dans V2, les fluctuations à chaque instant $t$ sont \textbf{indépendantes du passé} (marginales gaussiennes simples).
\end{remark}

%%%%%%%%%%%%%%%%%%%%%%%%%%%%%%%%%%%%%%%%%%%%%%%%%%%%%%%%%%%%%%%%%%%%%%%%%%%%%%%
\part{Implications pour les Garanties de Progrès}
%%%%%%%%%%%%%%%%%%%%%%%%%%%%%%%%%%%%%%%%%%%%%%%%%%%%%%%%%%%%%%%%%%%%%%%%%%%%%%%

\section{Lock-Free vs Wait-Free : Clarification}

\begin{definition}[Wait-Free vs Lock-Free]
\begin{itemize}
    \item \textbf{Wait-free} : Toute opération termine en un nombre borné d'étapes \textbf{avec succès garanti}.
    \item \textbf{Lock-free borné} : Toute opération termine en un nombre borné d'étapes, mais peut \textbf{retourner échec} (NULL si queue pleine/vide).
\end{itemize}
\end{definition}

\begin{remark}[Terminologie correcte]
L'implémentation MPMC est \textbf{lock-free bornée}, pas wait-free. Les opérations \texttt{reserve()} et \texttt{claim()} peuvent retourner NULL si la queue est pleine ou vide. Le succès dépend de l'état du système.
\end{remark}

\begin{theorem}[Producteur Lock-Free Borné dans V2]\label{thm:producer_wait_free}
L'opération \texttt{reserve()} du producteur est lock-free bornée avec complexité $O(S)$, où $S$ est le nombre de slots par partition.

\begin{proof}
L'algorithme V2 généralisé pour un producteur sur la partition $p$ avec $S$ slots est :
\begin{lstlisting}
reserve():
    for s = 0 to S-1:              // Borne: S iterations
        slot = &partition[p].slots[s]
        cur = READ(slot->state_gen)    // O(1)
        if state_of(cur) != EMPTY:
            continue                   // O(1)
        next = pack(gen, FILLING)
        if CAS(slot->state_gen, cur, next) == SUCCESS:
            return slot                // O(1)
    return FAIL (ou backpressure)
\end{lstlisting}

\textbf{Nombre d'étapes} : Au plus $3S$ opérations (READ, test, CAS) pour S slots.

\textbf{Boucle bornée} : La boucle est bornée à $S$ itérations maximum.

\textbf{Pas de contention producteur} : La partition $p$ n'est jamais touchée par un autre producteur.

\textbf{Dépendance sur l'état} : Le succès dépend de l'état laissé par les consommateurs. Si tous les S slots sont en état FILLING/READY/READING, l'opération retourne FAIL. Ce n'est pas une dépendance sur d'autres producteurs, mais sur la vitesse des consommateurs.

\textbf{Cas particulier $S=1$} : Retrouve la complexité $O(1)$ du modèle original.
\end{proof}
\end{theorem}

\begin{theorem}[Consommateur Lock-Free Borné dans V2]\label{thm:consumer_wait_free}
L'opération \texttt{claim()} du consommateur est lock-free bornée avec complexité $O(P \times S)$, où $P$ est le nombre de partitions et $S$ le nombre de slots par partition.

\begin{proof}
L'algorithme V2 généralisé pour un consommateur est :
\begin{lstlisting}
claim():
    for p = 0 to P-1:              // Borne: P partitions
        for s = 0 to S-1:          // Borne: S slots
            slot = &partition[p].slots[s]
            cur = READ(slot->state_gen)   // O(1)
            if state_of(cur) != READY:
                continue                  // O(1)
            next = pack(gen+1, READING)
            if CAS(slot->state_gen, cur, next) == SUCCESS:
                return (p, s)             // O(1)
            // Si CAS echoue, on passe au suivant (pas de retry!)
    return EMPTY
\end{lstlisting}

\textbf{Nombre d'étapes} : Au plus $3 \times P \times S$ opérations (READ, test, CAS).

\textbf{Boucle doublement bornée} : $P$ partitions $\times$ $S$ slots = $P \times S$ itérations max.

\textbf{Pas de retry} : Si le CAS échoue (un autre consommateur a pris le slot), on passe au suivant.

\textbf{Cas particulier $S=1$} : Retrouve la complexité $O(N)$ du modèle original (avec $N = P$).
\end{proof}
\end{theorem}

\section{Throughput Stationnaire}

\begin{theorem}[Throughput V2 --- Approximation]\label{thm:throughput_v2}
Sous l'hypothèse i.i.d., le throughput stationnaire du système V2 est approximativement :
\[
\boxed{\mathrm{Th}^* \approx N \cdot \lambda_C \cdot \pi_R^*}
\]
où $\pi_R^*$ est la probabilité stationnaire qu'un buffer soit dans l'état Ready.

\textbf{Limitation} : En pratique, $\lambda_C$ n'est pas uniforme --- les consommateurs scannent round-robin, donc un slot READY peut attendre si d'autres slots sont devant. Le throughput réel peut être inférieur.
\end{theorem}

\begin{proof}
À l'équilibre :
\begin{itemize}
    \item Chaque buffer est en état $R$ avec probabilité $\pi_R^*$
    \item En moyenne, $N \cdot \pi_R^*$ buffers sont en état $R$
    \item Chaque buffer $R$ est consommé avec taux $\lambda_C$
    \item Throughput total : $N \cdot \pi_R^* \cdot \lambda_C$
\end{itemize}
\end{proof}

\begin{corollary}[Scalabilité linéaire]
Sous les hypothèses du modèle, le throughput V2 est \textbf{linéaire en N} sans dégradation due à la contention producteur.

En V1, le throughput dégradait avec $T$ threads à cause de la contention sur l'état global.

\textbf{Note} : L'implémentation actuelle n'a pas d'état global partagé (\texttt{ready\_count}, \texttt{cons\_part} ont été supprimés). La seule contention consommateur vient des CAS sur les slots individuels.
\end{corollary}

%%%%%%%%%%%%%%%%%%%%%%%%%%%%%%%%%%%%%%%%%%%%%%%%%%%%%%%%%%%%%%%%%%%%%%%%%%%%%%%
\part{Récapitulatif}
%%%%%%%%%%%%%%%%%%%%%%%%%%%%%%%%%%%%%%%%%%%%%%%%%%%%%%%%%%%%%%%%%%%%%%%%%%%%%%%

\section{Comparaison des Preuves V1 vs V2}

\begin{table}[htbp]
\centering
\caption{Résumé comparatif des preuves.}
\small
\begin{tabular}{lcc}
\toprule
\textbf{Aspect} & \textbf{V1 (Mean-Field)} & \textbf{V2 (Indépendant)} \\
\midrule
Hypothèses & (H1)-(H4) & Indépendance seulement \\
Outil principal & Théorème de Kurtz & Loi des Grands Nombres \\
Martingale & $M^N(t)$, contrôle BDG & Non nécessaire \\
Convergence & $O(N^{-1/2})$ polynomiale & $O(e^{-cN})$ exponentielle \\
TCL & Ornstein-Uhlenbeck & Gaussien i.i.d. \\
ODE limite & $\dot{\rho} = F(\rho)$ & $\dot{\pi} = F(\pi)$ \\
Lock-free borné & Non garanti & \textbf{Garanti} \\
\bottomrule
\end{tabular}
\end{table}

\section{Message Clé}

\begin{center}
\fbox{\parbox{0.9\textwidth}{
\textbf{Le modèle V2 (No Producer Contention) est plus simple et plus puissant que V1.}

En éliminant la contention producteur :
\begin{enumerate}
    \item Les buffers producteur deviennent \textbf{indépendants}
    \item La preuve utilise la \textbf{LGN} au lieu du mean-field
    \item La convergence est \textbf{exponentiellement rapide} (sous i.i.d.)
    \item Les opérations sont \textbf{lock-free bornées par construction}
    \item Le throughput est \textbf{approximativement linéaire en N}
\end{enumerate}

\textbf{Limitations} : Le modèle suppose des taux de transition constants ($\lambda_C$), mais en pratique le scan round-robin peut créer des latences variables selon la position du slot. Le modèle i.i.d. est une approximation.
}}
\end{center}

\end{document}
